\documentclass[a4paper,11pt]{article}
\usepackage[utf8]{inputenc}
\usepackage[spanish]{babel}
\usepackage[margin=1cm]{geometry}
\usepackage{amsmath}
\usepackage{amsfonts}

\begin{document}
\pagestyle{empty}

\begin{minipage}[t][0.48\textheight][t]{0.48\textwidth}
  \centering \textbf{Evaluación Electrotecnia 2 - Tema A} \\
  \rule{\textwidth}{0.4pt}
  \begin{enumerate}
    \item Calcular la resistencia de un conductor con $\rho=1,7
      \times 10^{-8} \Omega m$, longitud $L=10m$ y sección $A=2
      \times 10^{-6} m^2$. (1 pt)
    \item Un circuito eléctrico está compuesto por una resistencia
      $R_1=10\Omega$ conectada en serie con un grupo de dos
      resistencias, $R_2=20\Omega$ y $R_3=20\Omega$, que se
      encuentran conectadas entre sí en paralelo. Determine resist.
      equiv. (2 pts)
    \item Se tiene una resistencia $R_1=5\Omega$ conectada en
      paralelo con un bloque formado por dos resistencias en serie,
      $R_2=10\Omega$ y $R_3=10\Omega$. Calcule la resistencia
      equivalente total. (2 pts)
    \item Calcular la inducción magnética $B$ en un solenoide de
      $N=500$ espiras, $I=2A$ y $l=0,2m$ ($\mu_0=4\pi \times 10^{-7}
      T\cdot m/A$). (1,5 pts)
    \item Calcular el flujo magnético $\Phi$ para un campo $B=1,5T$ y
      un área $A=0,01m^2$. (1 pt)
    \item Calcular la intensidad de campo $H$ para un solenoide de
      $N=1000$ espiras, $I=5A$ y $l=0,5m$. (1 pt)
    \item Si la intensidad de campo es $H=2000 Av/m$ y la
      permeabilidad relativa es $\mu_r=500$, calcular la inducción
      $B$. (1,5 pts)
  \end{enumerate}
\end{minipage}
\hfill
\begin{minipage}[t][0.48\textheight][t]{0.48\textwidth}
  \centering \textbf{Evaluación Electrotecnia 2 - Tema B} \\
  \rule{\textwidth}{0.4pt}
  \begin{enumerate}
    \item Calcular la inducción magnética $B$ en un solenoide de
      $N=600$ espiras, $I=1,5A$ y $l=0,3m$ ($\mu_0=4\pi \times
      10^{-7} T\cdot m/A$). (1,5 pts)
    \item Calcular la intensidad de campo $H$ para un solenoide de
      $N=800$ espiras, $I=4A$ y $l=0,4m$. (1 pt)
    \item Calcular el flujo magnético $\Phi$ para un campo $B=0,8T$ y
      un área $A=0,02m^2$. (1 pt)
    \item Se tiene una resistencia $R_1=15\Omega$ conectada en serie
      con un grupo de dos resistencias en paralelo, $R_2=15\Omega$ y
      $R_3=15\Omega$. Calcule la resistencia equivalente total del
      circuito. (2 pts)
    \item Un circuito presenta una resistencia $R_1=10\Omega$ en
      paralelo con un bloque de dos resistencias en serie,
      $R_2=10\Omega$ y $R_3=10\Omega$. Determine la resistencia total. (2 pts)
    \item Calcular la resistencia de un conductor con $\rho=1,7
      \times 10^{-8} \Omega m$, longitud $L=20m$ y sección $A=1
      \times 10^{-6} m^2$. (1 pt)
    \item Si la intensidad de campo es $H=1500 Av/m$ y la
      permeabilidad relativa es $\mu_r=800$, calcular la inducción
      $B$. (1,5 pts)
  \end{enumerate}
\end{minipage}

\vfill

\begin{minipage}[t][0.48\textheight][t]{0.48\textwidth}
  \centering \textbf{Evaluación Electrotecnia 2 - Tema C} \\
  \rule{\textwidth}{0.4pt}
  \begin{enumerate}
    \item En un montaje eléctrico, se conectan dos resistencias de
      $30\Omega$ y $60\Omega$ en paralelo. El resultado de esta
      conexión se coloca en serie con una resistencia de $10\Omega$.
      Calcule la resistencia total. (2 pts)
    \item Calcular el flujo magnético $\Phi$ para un campo $B=2,0T$ y
      un área $A=0,005m^2$. (1 pt)
    \item Calcular la resistencia de un conductor con $\rho=1,7
      \times 10^{-8} \Omega m$, longitud $L=5m$ y sección $A=4 \times
      10^{-6} m^2$. (1 pt)
    \item Si la intensidad de campo es $H=3000 Av/m$ y la
      permeabilidad relativa es $\mu_r=400$, calcular la inducción
      $B$. (1,5 pts)
    \item Se tiene una resistencia $R_1=20\Omega$ conectada en
      paralelo con un bloque formado por dos resistencias en serie,
      $R_2=20\Omega$ y $R_3=20\Omega$. Calcule la resistencia
      equivalente total. (2 pts)
    \item Calcular la inducción magnética $B$ en un solenoide de
      $N=400$ espiras, $I=3A$ y $l=0,1m$ ($\mu_0=4\pi \times 10^{-7}
      T\cdot m/A$). (1,5 pts)
    \item Calcular la intensidad de campo $H$ para un solenoide de
      $N=1200$ espiras, $I=3A$ y $l=0,6m$. (1 pt)
  \end{enumerate}
\end{minipage}
\hfill
\begin{minipage}[t][0.48\textheight][t]{0.48\textwidth}
  \centering \textbf{Evaluación Electrotecnia 2 - Tema D} \\
  \rule{\textwidth}{0.4pt}
  \begin{enumerate}
    \item Calcular la inducción magnética $B$ en un solenoide de
      $N=700$ espiras, $I=1A$ y $l=0,25m$ ($\mu_0=4\pi \times 10^{-7}
      T\cdot m/A$). (1,5 pts)
    \item Calcular la resistencia de un conductor con $\rho=1,7
      \times 10^{-8} \Omega m$, longitud $L=15m$ y sección $A=3
      \times 10^{-6} m^2$. (1 pt)
    \item Calcular la intensidad de campo $H$ para un solenoide de
      $N=600$ espiras, $I=6A$ y $l=0,3m$. (1 pt)
    \item Un circuito eléctrico está compuesto por una resistencia
      $R_1=5\Omega$ conectada en serie con un grupo de dos
      resistencias en paralelo, $R_2=10\Omega$ y $R_3=10\Omega$.
      Determine la resistencia total. (2 pts)
    \item Si la intensidad de campo es $H=1000 Av/m$ y la
      permeabilidad relativa es $\mu_r=1000$, calcular la inducción
      $B$. (1,5 pts)
    \item Se tiene una resistencia $R_1=10\Omega$ conectada en
      paralelo con un bloque formado por dos resistencias en serie,
      $R_2=5\Omega$ y $R_3=5\Omega$. Calcule la resistencia
      equivalente total. (2 pts)
    \item Calcular el flujo magnético $\Phi$ para un campo $B=1,2T$ y
      un área $A=0,015m^2$. (1 pt)
  \end{enumerate}
\end{minipage}

\end{document}

